\scanpage{11}

$X$ surface topologique

$\mathcal{G}$ groupe discret opérant sur $X$

$G$ sous-groupe invariant de $\mathcal{G}$, opérant proprement sur $X$, de sorte que
\unsure{\deleted{$\forall x\in X$}}
\deleted{\illegible{1}} les stabilisateurs $G_x$ $(x\in X)$ \unsure{opèrent} fidèlement au voisinage de $x$, en conservant l'orientation locale en $x$

$Y=X/G$ (\unsure{surface topologique}), $f:X\to Y$

$X^{!}$ partie discrète de $X$, stable sous $\mathcal{G}$,

$d_X$ ``signature'' sur $X$ de support $X^{!}$, inv. par $\mathcal{G}$
\[
\unsure{a\mapsto d_X(a)}\qquad X^{!}\longmapsto\mathbb{N}
\]

$Y^{!}\simeq X^{!}/G\subset Y$ image de $X^{!}$ dans $Y$

On suppose $x\in X$, $G_x\ne\{1\}\Rightarrow x\in X^{!}$

\[
d_Y:Y^{!}\to\mathbb{N},\qquad d_Y(y)=\nu_{\unsure{a|x}}d_X(x)
\]
si $y=f(x)$, $\nu_{\unsure{a|x}}=[G_x]$

$\operatorname{Rev}(X,d_X,\mathcal{G})$ catégorie des $\mathcal{G}$-revêtements ramifiés de $X$, subordonnés à $d_X$

$\operatorname{Rev}(Y,d_Y,\mathcal{H})$ catégorie des \unsure{$\mathcal{H}$-rev\textsuperscript{ts}} ramifiés de $Y$, subordonnés à $d_Y$

Pour $X'\in\operatorname{Ob}\operatorname{Rev}(X,d_X,\mathcal{G})$, considérons $X'/G$ comme surface sur laquelle opère $\mathcal{H}=\mathcal{G}/G$, au dessus de $Y=X/G$. Alors

\margin{10}

\scanpage{12}
\mspage{294}

a) $X'/G\in\operatorname{Ob}\operatorname{Rev}(Y,d_Y,\mathcal{H})$

b) On \unsure{trouve} ainsi un foncteur
\[
(19)\qquad\operatorname{Rev}(X,d_X,\mathcal{G})\longrightarrow\operatorname{Rev}(Y,d_Y,\mathcal{H})
\]

c) Ce foncteur est une équivalence de catégories.

d) Description d'un foncteur quasi-inverse par image inverse ``normalisée''.

Remarques. \unsure{OPS|OFS} $d_X$ à valeurs $>0$ i.e. $d_Y$ à valeurs $>0$ (sinon, si $S=\{x\in X^{!}\mid d_X(x)=0\}$, \unsure{remplaçons} $X$ par $X\setminus S$), $Y$ par $Y\setminus T$ $(T=S/G)$.

\inserted{\unsure{OPS|OFS} $d_X(x)=1\Rightarrow[G_x]\ne1$ (sinon, remplacer $X^{!}$ \unsure{par} $X^{!}\setminus S$
\unsure{$S=\{x\in X^{!}\mid d_X(x)=1,[G_x]=1$}}

Le cas le plus intéressant \unsure{pour nous est} celui où \deleted{$d_X(x)=1$} $\forall x\in X^{!}$, on a $d_X(x)=1$ et $G_x\ne\{1\}$ (donc $X^{!}=\{x\in X\mid G_x\ne\{1\}\}$), dans le cas où
\[
(20)\qquad\operatorname{Rev}(X,d_X,\mathcal{G})=\operatorname{Rev}(X,\mathcal{G})
\]
et où $\operatorname{Rev}(Y,d_Y,\mathcal{H})$ peut donc s'interpréter comme un $\operatorname{Rev}(X,\mathcal{G})$ (sans \unsure{donnée} de ramification)

\margin{11}

\scanpage{13}

\uline{Corollaire} Soit $x\in X\setminus X^{!}$, $y$ son image dans $Y\setminus Y^{!}$, on a un iso. canonique
\[
(21)\qquad
\Pi_1((X,d_X,\mathcal{G})\unsure{,|;}x)\xrightarrow{\sim}\Pi_1((Y,d_Y,\mathcal{H})\unsure{,|;}y)
\]

Nous allons expliciter cet isomorphisme en termes de \unsure{calculs} de \deleted{lacets}\inserted{chemins}.

Un élément \inserted{$\alpha$} de $\Pi_1((X,d_X,\mathcal{G})\unsure{,|;}x)$ \unsure{\deleted{provient}} d'un couple $(\gamma,\ell)$, avec
\[
\gamma\in\mathcal{G},\qquad \ell\text{ classe de chemins dans }X\setminus X^{!}\text{ de }x\text{ vers }\gamma^{-1}x
\]

(avec $(\gamma,\ell)$ et $(\gamma',\ell')$ définissant le $\overline{\mathrm m}$ élément \unsure{sss|ssi} $\gamma'=\gamma$, et \inserted{l'élément} ${\ell'}^{-1}\ell$ de $\Pi_1(X\setminus X^{!},x)$ est dans le ss-groupe invariant noyau de
\[
\Pi_1(X\setminus X^{!},x)\longrightarrow\Pi_1((X,d_X),x)
\]
--- i.e le ss-groupe invariant engendré par les \unsure{\uline{le}|\uline{$\ell$}} $\lambda_a^{d_X(a)}$, où pour $\forall a\in X^{!}$, $\lambda_a$ désigne la \unsure{classe} dans $\Pi_1(X\setminus X^{!},x)$ d'un ``lacet de $x$ autour de $a$'').

Soit alors $\dot\gamma$ l'image de $\mathcal{G}$ dans $\mathcal{H}$ --- \unsure{Et|et} soit
\[
f_U:X\setminus X^{!}\overset{\unsure{\mathrm{dfn}|\mathrm{d\acute{e}f}}}{=}U\longrightarrow Y\setminus Y^{!}\overset{\unsure{\mathrm{dfn}|\mathrm{d\acute{e}f}}}{=}V\simeq U/G
\]
l'application

\margin{12}

\scanpage{14}
\mspage{295}

canonique, alors l'image \inserted{$\beta$} de \inserted{$\alpha\in\Pi_1(X,d_X,\mathcal G;x)$ défini par} $(\gamma,\ell)$ \deleted{\illegible{1}} dans $\Pi_1(Y,d_Y,\mathcal H;y)$ est $(\dot\gamma,f_U(\ell))$

(\uline{NB} $f_U(\ell)$ est une classe de chemins dans $V=Y\setminus Y^{!}$ de $y=f_U(x)$ vers $\dot\gamma^{-1}(y)=f_U(\gamma^{-1}x)$).

\margin{\uline{NB.}}
On voit tout de suite que si $(\gamma,\ell)$ et $(\gamma',\ell')$ définissent un $\overline{\mathrm m}$ $\alpha$, alors ils définissent aussi un $\overline{\mathrm m}$ $\beta$, ce qui provient du fait que
\[
\Pi_1(f_U,x):\Pi_1(X\setminus X^{!},x)\longrightarrow\Pi_1(Y\setminus Y^{!},y)
\]
applique un $\lambda_a$ dans \inserted{$\unsure{\rightsquigarrow}$} $\lambda_b^{\nu(a)}$
($b=f(a)$, $\nu(a)=[G_a]={}$degré de ramification de $f$ en $a$) donc
\[
\lambda_a^{d_X(a)}\ \unsure{\rightsquigarrow|\sim}\ \lambda_b^{d_X(a)\nu(a)=d_Y(b)}.
\]

\unsure{Revenons} à la situation de (13)
\[
X=U_{0,3}^{\mathrm{an}},\quad
\mathcal G=\mathcal T_{0,3}=\unsure{\mathfrak S_3|\mathcal S_3}\times\mathbb Z/2,\quad
G=\mathcal T_{0,3}^{+}=\unsure{\mathfrak S_3|\mathcal S_3}
\]
\inserted{$X^{!}=\{\deleted{\illegible{1}}\,2,-1,\frac12;-\bar j,-j\}$}
$Y\simeq\mathbb P^{1\,\mathrm{an}}\setminus\{0\}$, $Y^{!}=\{1,\infty\}$, \inserted{$\mathcal H=\mathbb Z/2\mathbb Z$,} application de \unsure{passage} au quotient de $\widehat X=\mathbb P^{1\,\mathrm{an}}$ dans $\widehat Y=\mathbb P^{1\,\mathrm{an}}$ donnée par (13)
\[
f(z)=\frac{3^3}{2^2}\frac{z^2(z-1)^2}{(z^2-z+1)^3}.
\]
\margin{13}

\scanpage{15}

Pour calculer $\Pi_1(Y,\unsure{\underline d_Y|d_Y},\mathbb Z/2\mathbb Z)=\widetilde\Pi'_{0,3}$ (groupe \unsure{cartographique} triangulé non orienté), on a pris le point base $y=-\bar j$, on va donc choisir un $x\in U_{0,3}^{\mathrm{an}}$ au dessus de $y$, i.e. tel que
\[
f(x)=y=-\bar j
\]
--- il y a \deleted{six \unsure{places|points}} \inserted{un} choix dans un \unsure{ens} $f^{-1}(\{y\})$ à six éléments, dont \uline{trois} dans l'hémisphère nord (les trois autres dans l'hémisphère sud), ceux-ci sont permutés transitivement par $\rho$.
Ils se trouvent \unsure{respectivement} dans les triangles sphériques $(-\bar j,0,\frac12)$, $(-\bar j,1,2)$, $(-\bar j,\infty,-1)$ (appliqués par $f$ homéomorphiquement sur le triangle sphérique $(0,1,\infty)$ \unsure{supérieur}, \inserted{(ou l'\unsure{hémisphère|hémisph} sup)} de centre $-\bar j$). Le choix de $x$ revient donc au choix de l'un des trois \inserted{triangles sphériques, ou encore des $3$} segments sphériques $(0,\frac12)$, $(1,2)$, $(\infty,-1)$ --- j'ai envie de prendre $(0,\frac12)$, \unsure{correspondant} au triangle $(-\bar j,0,\frac12)$.

Les trois générateurs \inserted{involutifs} \deleted{\illegible{1}} $\tau'_0,\tau'_1,\tau'_\infty$ de $\Pi_1(Y,d_Y,\mathbb Z/2\mathbb Z)$ \inserted{\unsure{au dessus} de $\tau\in\mathbb Z/2\mathbb Z$} s'obtiennent en prenant des chemins $\ell'_0,\ell'_1,\ell'_\infty$ du \inserted{pôle nord} $y=-\bar j$ vers \inserted{le pôle sud} $\tau y=-j$, qui sont les arcs de \unsure{grand cercle|grands cercles} coupant l'équateur resp. en $2,-1,\frac12$.
\margin{14}

\scanpage{16}
\mspage{296}

Ces chemins sont les images par $f$ de chemins sur $X$, soient $\lambda_0:\overset{\curvearrowright}{x,z_0}$, \inserted{$\lambda_1:$} $\overset{\curvearrowright}{x,z_1}$, \inserted{$\lambda_\infty:$} $\overset{\curvearrowright}{x,z_\infty}$, qui coupent \inserted{chacun} la subdivision barycentrique \unsure{du double|des doubles} \unsure{triangle sphérique|triangles sphériques} $(0,1,\infty)$ en exactement un point, situé \unsure{respectivement} sur \unsure{le côté|les côtés} $(\frac12,-\bar j)$, $(-\bar j,0)$ et $(0,\frac12)$ du triangle sphérique $(0,\frac12,-\bar j)$.

\margin{\uline{NB} les points $z_0,z_1,z_\infty$ sont \unsure{caractérisés} \inserted{\unsure{resp.}} comme les seuls points de $X$ au dessus de $-j=f(\dot\tau x)$, \unsure{se trouvant} dans les triangles sphériques $(-\bar j,\frac12,1)$, $(0,-\bar j,-1)$, $(\frac12,0,-j)$.}

Les points \deleted{\illegible{1}} $z_0,z_1,z_\infty$ au dessus de $\tau y$ sont \unsure{respectivement} de la forme \inserted{$g_0x={}$} $\dot\sigma_\infty\dot\tau x$, \inserted{$g_1x={}$} $\dot\sigma_0\dot\tau x$, \inserted{$g_\infty x={}$} \deleted{$\unsure{\dot\sigma_\infty}$} $\dot\tau x$.
\inserted{(i.e. $g_0=\dot\sigma_\infty\dot\tau$, $g_1=\dot\sigma_0\dot\tau$, $g_\infty=\dot\tau$)}
Ainsi les générateurs involutifs $\tau'_0,\tau'_1,\tau'_\infty$ de $\widetilde\Pi'_{0,3}=\Pi_1(Y,d_Y,\mathbb Z/2\mathbb Z;y)$ sont les images des éléments \deleted{$(\unsure{\dot\sigma_0\dot\tau},\unsure{\lambda_0})$}, (notés par le $\overline{\mathrm m}$ signe)
\[
\tau'_0=(\dot\sigma_0\dot\tau,\lambda_0);\qquad
\tau'_1=(\dot\sigma_1\dot\tau,\lambda_1),\qquad
\tau'_\infty=(\dot\sigma_\infty\dot\tau,\lambda_\infty)
\]
de $\Pi_1(X,\mathcal G;x)$. Nous allons utiliser l'arc de grand cercle \inserted{$\ell$} de \deleted{\unsure{$x$ à}} $-\bar j=P$ (\deleted{origine} \inserted{point-base} habituel pour le calcul de $\Pi_1(X,\mathcal G)$) \inserted{à $x$,} pour identifier $\Pi_1(X,\mathcal G;x):\Pi_1(X,\mathcal G;\unsure{P|p})$.

On peut alors prendre \deleted{\illegible{1}}
\[
\tag{21}
\tau'_0=(\dot\sigma_0\dot\tau,\lambda'_0),\qquad
\tau'_1=(\dot\sigma_1\dot\tau,\lambda'_1),\qquad
\tau'_\infty=(\dot\sigma_\infty\dot\tau,\lambda'_\infty)
\]
où $\lambda'_0,\lambda'_1,\lambda'_\infty$ sont les \deleted{\illegible{2} arcs de grand cercle (\unsure{avec les bouts|les deux bouts} en $P$)} \inserted{chemins \unsure{de}} de $P=-\bar j$ \inserted{vers \deleted{\illegible{1}} $g_iP$ \deleted{\illegible{1}}} \deleted{vers $z_0,z_1,z_\infty$ \unsure{respectivement} (qui, à l'excep-}
\margin{15}

\scanpage{17}

\deleted{tion \unsure{du dernier|des derniers}, ne \unsure{rencontrent} pas le \unsure{graphe} barycentrique \unsure{du double triangle sphérique|des doubles triangles sphériques} sauf en $P=-\bar j$, le dernier le \unsure{rencontrant} \inserted{\unsure{seult}} en un seul autre point, étant sur le segment sphérique $(0,\frac12)$. Les points $z_0,z_1,z_\infty$ sont caractérisés respectivement comme les points $z$ tels que $f(z)=\unsure{-j|-\bar j}$, et qui se trouvent sur les triangles sphériques $(\unsure{\frac12},\unsure{-\bar j},1)$, $(-\bar j,0,-1)$, $(0,\frac12,-j)$.}

\figure{[Deux sphères dessinées l'une au-dessus de l'autre, séparées par une flèche verticale descendante portant $f$. Les pôles sont marqués $-\bar j$ en haut et $-j$ en bas. Sur la sphère supérieure, des grands cercles et arcs en traits pleins et pointillés passent par les points marqués $0,1,\infty,-1,\frac12,2$; le point $x$ est à l'intérieur d'un triangle au-dessus de l'équateur, $z_0$ à droite, $z_1$ à gauche et $z_\infty$ au-dessous de l'équateur. Des chemins joignent $x$ à ces trois points, avec les étiquettes \unsure{$\lambda_0,\lambda_1,\lambda_\infty$}; un chemin joignant $-\bar j$ à $x$ porte \unsure{$\ell$}. Autres petites étiquettes près des arcs supérieurs: \unsure{$\lambda'_1,\lambda'_\infty$}, \illegible{2}. La sphère inférieure porte les mêmes six marques équatoriales et des méridiens pointillés; un arc appuyé va de $0$ vers $1$, et un trait épais recourbé est tracé près de $1$.]}

\unsure{donnés} par
\[
\lambda'_i=(\deleted{\illegible{1}}\unsure{g_i}(\unsure{\ell^{-1}|\ell'}))(\lambda_i)\ell.
\]

On a donc
\[
\lambda'_0=
\underbrace{\dot\sigma_\infty\dot\tau}_{\substack{\text{symétrie}\\\text{\unsure{p.r. à \uline{l'équateur}}}\\(\infty,-\bar j,\frac12,-j)}}
(\ell^{-1})\circ\lambda_0\circ\ell
\qquad
\begin{gathered}
\deleted{\illegible{1}}\\
\text{homotope au}\\
\text{lacet trivial}\\
\text{dans }U_{0,3}
\end{gathered}
\]
\[
\lambda'_1=
\underbrace{\dot\sigma_0\dot\tau}_{\substack{\text{symétrie}\\\text{\unsure{p.r. à} l'équateur}\\(0,-\bar j,2,-j)}}
(\ell^{-1})\circ\lambda_1\circ\ell
\qquad
\begin{gathered}
\text{homotope au}\\
\text{lacet trivial}\\
\text{dans }U_{0,3}
\end{gathered}
\]
\[
\lambda'_\infty=(\underbrace{\tau}(\ell^{-1}))\circ\lambda_\infty\circ\ell
\qquad
\begin{gathered}
\text{homotope au}\\
\deleted{\illegible{2}}\inserted{\text{arc de}}\text{ grand}\\
\text{cercle }(-\bar j,\frac12,-j)\\
\text{de }P=-\bar j\text{ vers}\\
\dot\tau P=-j
\end{gathered}
\]

\margin{16}

\scanpage{18}
\mspage{297}

En résumé, avec les notations de [C] \deleted{\unsure{(13)|(43)}} p. 7 \unsure{ff|H}, on aura donc
\[
(22)\quad
\left\{
\begin{aligned}
\tau'_0&=\tilde\sigma_\infty &&(\text{sur }\dot\sigma_\infty\dot\tau)\\
\tau'_1&=\tilde\sigma_0 &&(\text{sur }\dot\sigma_0\dot\tau)\\
\tau'_\infty&=\tau_\infty &&(\text{sur }\dot\tau)
\end{aligned}
\right.
\]
\margin{NB $\tau_\infty=\sigma_\infty\tilde\sigma_\infty=\tilde\sigma_\infty\sigma_\infty$.}

Il faut vérifier les relations essentielles
\[
(23)\quad
{\tau'_0}^{2}={\tau'_1}^{2}={\tau'_\infty}^{2}
=\bigl(\underbrace{\tau'_0\tau'_\infty}_{\sigma_\infty}\bigr)^2
=\bigl(\underbrace{\tau'_1\tau'_0}_{\rho^{-1}}\bigr)^3=1
\]
les quatre premières sont claires, la dernière s'écrit
\[
(\tilde\sigma_0\tilde\sigma_\infty)^3=1
\]
\unsure{or on a}
\[
\tilde\sigma_0\tilde\sigma_\infty=\rho^{-1}\qquad
(\inserted{\text{[C],}}\text{ formule }(13))
\]
d'où ce qu'on veut.

Calculons (conformément à (18)) $\tau'_\infty\tau'_1,\tau'_0\tau'_\infty,
\underbrace{\tau'_1\tau'_0}_{\rho^{-1}}$ --- qui devront correspondre à $\ell'_0,\ell'_1,\ell'_\infty$ ---
\[
\begin{aligned}
\ell'_0&=\tau'_\infty\tau'_1=\tau_\infty\tilde\sigma_0
=\sigma_\infty\underbrace{\tilde\sigma_\infty\tilde\sigma_0}_{\rho}
=\sigma_\infty\rho\\
\ell'_1&=\tau'_0\tau'_\infty=\unsure{\sigma_\infty}\\
\ell'_\infty&=\tau'_1\tau'_0=\rho^{-1}
\end{aligned}
\]
i.e.
\[
(24)\quad \ell'_0=\sigma_\infty\rho,\quad \ell'_1=\sigma_\infty,\quad \ell'_\infty=\rho^{-1}
\]
avec les relations caractéristiques
\[
(25)\quad \ell'_\infty\ell'_1\ell'_0=1,\qquad
{\ell'_{\unsure{1|\infty}}}^{2}=1,\qquad {\ell'_\infty}^{3}=1
\]
Tout est OK !

\margin{On exprime les éléments remarquables $\rho,\sigma_0,\sigma_1,\sigma_\infty,\ell_0,\ell_1,\ell_\infty,\tau_0,\tau_1,\tau_\infty$ de $\Pi_1(U_{0,3}^{\unsure{\mathrm{an}}};\unsure{\mathfrak S_3}\times\unsure{\mathbf Z|\mathbb Z}/2)$ en fonction des $\tau'_0,\tau'_1,\tau'_\infty$ ainsi :
\[
\rho=\tau'_0\tau'_1
\]
\[
\tau_\infty=\tau'_\infty,\quad
\tau_0=\rho\tau_\infty\rho^{-1}=\tau'_0\tau'_1\tau'_\infty\tau'_1\tau'_0,
\]
\[
\tau_1=\rho^{-1}\tau_\infty\rho=\ldots=\tau'_1\tau'_\infty\tau'_1,
\]
\[
\sigma_\infty=\tau'_0\tau'_\infty\inserted{(=\tau'_\infty\tau'_0)},\quad
\sigma_0=\rho\sigma_\infty\rho^{-1}=\tau'_0\tau'_1\tau'_0\tau'_\infty\unsure{\tau'_1\tau'_0}
\]
\[
\sigma_1=\rho^{-1}\sigma_\infty\rho=\tau'_1\tau'_\infty\tau'_0\tau'_1
\]
\[
\ell_0=\tau_\infty\tau_1=\tau'_\infty\tau'_1\tau'_\infty\tau'_1=(\tau'_\infty\tau'_1)^2
\]
\[
\ell_1=\tau_0\tau_\infty=\tau'_0\tau'_1\tau'_\infty\tau'_1\tau'_\infty\unsure{\tau'_0}
\]
\[
\ell_\infty=\tau_1\tau_0=(\underbrace{\tau'_1\tau'_\infty\tau'_1}\tau'_0)^2
\]
\[
\varepsilon_0=\ell'_0=\deleted{\illegible{3}}\inserted{\unsure{\tau'_\infty\tau'_1}}
\]
\[
\varepsilon_1=\tau'_0\tau'_1\tau'_\infty\tau'_0
\]
\[
\varepsilon_\infty=\tau'_1\tau'_\infty\tau'_1\tau'_{\illegible{1}}
\]
}

\margin{17}

\scanpage{19}

Je voudrais \unsure{encore} expliciter \unsure{les relations|la relation} de
$\widetilde{\Pi}'_{0,3}$
\inserted{\unsure{(groupe} carto-\unsure{graphique} non orienté)}
$\simeq\unsure{\mathfrak{G}_{0,3}}$
\inserted{$=\mathrm{Aut}_{\unsure{\mathrm{loc}|\mathrm{ext}}}(\Pi_{0,3})$}
avec
$\mathrm{GP}(1,\mathbb Z)\simeq\mathrm{Gl}(\unsure{2|1},\mathbb Z)'\mathrel{\mathop{=}\limits^{\mathrm{d\acute{e}f}}}\mathrm{Gl}(2,\mathbb Z)/\{\pm1\}$.
\unsure{Rappelons} d'abord l'opération de
\[
(26)\qquad
\mathrm{GP}(1,\mathbb R)\simeq\mathrm{Gl}(2,\mathbb R)/\mathbb R^*
\simeq\mathrm{Gl}(2,\mathbb R)^{\pm1}/\{\pm1\}
\qquad\deleted{\unsure{agit}}
\]
\margin{$\unsure{f}$}
sur le demi-plan de \unsure{Poincaré} $D$
\[
(27)\qquad D=\{z\in\mathbb C\mid\Im(z)>0\}
=\{x+iy\mid x\in\mathbb R,\;y\in\mathbb R^{*+}\}
\]
On fait opérer la matrice
$u=\left(\begin{array}{cc}a&b\\c&d\end{array}\right)$ par
\[
(28)\qquad f_u(z)=
\left\{\begin{array}{ll}
\dfrac{az+b}{cz+d}&\text{si }\det u>0\\[6pt]
\dfrac{a\bar z+b}{c\bar z+d}&\text{si }\det u<0
\end{array}\right.
\]
De cette façon, $\mathrm{GP}(1,\mathbb R)$ est isomorphe au groupe des automorphismes
\uline{conformes} de $D$.

\margin{\uline{NB} $\mathrm{GP}(1,\mathbb R)$ a deux \unsure{comp.} connexes,
la \unsure{comp.} neutre \inserted{$\mathrm{GP}(1,\mathbb R)^+$}
\unsure{est} l'image de $\mathrm{Sl}(2,\mathbb R)$
(\unsure{ou} de $\mathrm{Gl}(2,\mathbb R)^+$)
et s'identifie à $\mathrm{Sl}(2,\mathbb R)/\{\pm1\}$.}

D'ailleurs, la donnée d'une courbe elliptique \inserted{$E$} sur $\mathbb C$
équivaut à celle d'un espace $V$ de dim $1$ sur $\mathbb C$, et d'un
``réseau'' $\pi\subset V$ (sous-groupe discret de rg $2$).
\margin{\unsure{équivalence|équivalences} de catégories \illegible{1}.}
L'orientation canonique de $V$ \unsure{correspond} à un générateur \inserted{$\unsure{\eta}$} de
$\bigwedge_{\mathbb Z}^{2}\pi$ -- ainsi $\pi$ est ``orienté''.
Le choix d'une rigidification

\margin{18}

\scanpage{20}
\mspage{298}

de \unsure{type} unimodulaire de $\unsure{\pi|\Pi}$, équivaut à celui d'une base
$(e_1,e_2)$ de $\unsure{\pi|\Pi}$ \inserted{(comme $\mathbb Z$-module)} telle que
$e_1\wedge e_2=\unsure{\eta}$.
L'élément $e_1$ \inserted{(\unsure{comme base} de $V$ sur $\mathbb C$)}
permet d'ailleurs d'identifier $V$ à $\mathbb C$, $e_1$ devenant $1$.
\unsure{Ainsi,|Ainsi;} la donnée de $E$ avec sa rigidification de \unsure{type}
unimodulaire, équivaut à celle d'un élément $e_2=\tau$ de $\unsure{\mathbb C|E}$,
tel que

$1^\circ$) $(1,\tau)$ \unsure{engendrent|engendre} un sous-groupe discret de \unsure{rg} $1$,
i.e. $\tau\notin\mathbb R$ i.e. $\Im\tau\ne0$

$2^\circ$) l'orientation dudit sous-groupe définie par la base $(1,\tau)$ est
``\unsure{directe}'', i.e. $\Im\tau>0$, i.e. $\tau\in D$.

Quand on regarde l'opération de $\mathrm{Sl}(2,\mathbb Z)$ sur \inserted{\unsure{l'ens}}
des classes d'isomorphie des courbes ell. rigidifiées, \unsure{ens.} identifié à $D$,
on trouve \unsure{que|pour} $\left(\begin{array}{cc}a,&b\\c,&d\end{array}\right)$ opère par
$\tau\longmapsto\dfrac{a\tau+b}{c\tau+d}$, i.e. c'est l'opération habituelle de
$\mathrm{Sl}(2,\mathbb Z)$ sur $D$, à travers
$\mathrm{Sl}(2,\mathbb Z)'\simeq\mathrm{GP}(1,\mathbb Z)^+$.
Ce groupe opère proprement, Il a exactement deux trajectoires critiques, qui \unsure{corres-}
\margin{19}
